\documentclass{article}
% Competition/workshop paper, non-anonymous (single-blind) camera-ready: the official style
% requires a track option together with [final]; [final] alone leaves the track name undefined.
\usepackage[sglblindworkshop,final]{neurips_2026}
\workshoptitle{Smart Buildings Challenge (SB26)}
\usepackage[utf8]{inputenc}
\usepackage[T1]{fontenc}
\usepackage{url}
\usepackage{booktabs}
\usepackage{amsfonts}
\usepackage{amsmath}
\usepackage{nicefrac}
\usepackage{microtype}
\usepackage{xcolor}
\usepackage{graphicx}
\usepackage{algorithm}
\usepackage{algpseudocode}
\usepackage{multirow}
\usepackage{hyperref}  % after float/algorithm so caption anchors are not duplicated
\input{results_macros}
% Figure include that degrades to a visible placeholder when the file is not built yet.
\newcommand{\figfile}[2]{\IfFileExists{#2}{\includegraphics[width=#1]{#2}}{\fbox{\parbox{#1}{\centering\textbf{??} missing figure \texttt{\detokenize{#2}}}}}}

\title{Schedule-Aware Demand-Based Reset: A Structured Rule-Based Controller for Cross-Building HVAC Transfer in the Smart Buildings Challenge}

\author{%
  Neel Sirivara\\
  University of Washington\\
  \texttt{nsiriv@cs.washington.edu}
}

\begin{document}

\maketitle

\begin{abstract}
The NeurIPS 2026 Smart Buildings Challenge asks for one HVAC control policy that operates
procedurally generated commercial buildings it has never seen, under 30-day weather strips
from five climates, trading occupant comfort against energy cost and carbon. We present
SADR (Schedule-Aware Demand-based Reset), a pure-numpy rule-based controller that reads the
plant topology from the building metadata and is otherwise built from facts shared by every
generated building: comfort is charged only while zones are occupied, fan power grows as
static pressure to the three-halves, the chiller has a constant coefficient of performance,
and all heat reaches a zone through a reheat coil whose discharge is capped at 308.15\,K.
SADR gates its mode on occupancy and schedule (with a self-tuning optimum start), regulates each zone with a PI
loop on \emph{net heat} translated into damper and valve commands through the supply-air and discharge physics,
and resets pressures and hot-water temperature from zone demand. On the 30 public scenarios the round-4 version
(v7) scores \RFConfPubCand{} on the leaderboard metric (mean absolute per-step reward, lower is better) against
\MetricPidWater{} for the strongest organizer baseline (\texttt{pid\_water}), a paired-bootstrap ratio of
\RFPwPubRatio{} (95\% CI \RFPwPubRatioLo--\RFPwPubRatioHi); on \NHeldAll{} scenarios from \NHeldAllBuild{}
generated buildings the ratio is \RFPwOldRatio. Loss decompositions drove the revisions: coldest supply air and a
deeper setback (\CandWinnerRelRed{} on the public set); zone loops scaled with floor area for large zones
(\HeldAllVFiveRelRed{} on held-out buildings); and a pre-registered screen of discharge headroom, pressure-aware
requests, night air and set-back ventilation (\RFConfPoolRed{}, CI \RFConfPoolRedLo{} to \RFConfPoolRedHi, on
\RFConfPoolN{} held-out scenarios no decision touched). On Kaggle's hidden scenarios v7 scores \RFiveKagVSeven,
\RFiveKagGapPub{} times its public value; submissions of fixed earlier versions and of the baseline show that the
hidden buildings behave like the public ones and the hidden weather is harsher. A fifth round's deeper weekend
setback and lifted request cap (v8) lower the loss by \RFiveConfPubRelAbs{} on the public set (Kaggle:
\RFiveKagVEight). A sixth round on new harsh-weather strips, which together with the public ones reproduce the Kaggle
scores, ships v9 (cheaper heat for set-back zones through air the handlers move anyway, a deeper weeknight setback):
\RSixConfHarshCRelAbs{} lower on \RSixConfHarshCN{} unseen harsh scenarios (CI \RSixConfHarshCRelHiAbs{} to
\RSixConfHarshCRelLoAbs), a tie on the public set. We release the controller and report where it still loses.
\end{abstract}

\section{Introduction}

Learned building controllers have been studied for two decades, yet few run in real
buildings; the recurring obstacles are data-hungry training, fragile robustness and poor
generalization across buildings~\citep{wang2020rl}. The Smart Buildings Challenge (SB26)
isolates the last of these: a submission is a single policy class, instantiated on held-out
buildings drawn from the same procedural generator as six public ones and scored on 30-day
episodes under weather strips it has not seen~\citep{sb26kit}. The observation and action
vectors change shape with the building ($K$ air handlers and $V$ zones expose $K{+}1$ on/off
commands and $3K{+}2{+}2V$ continuous setpoints), and the simulator is Google's open-source
Smart Buildings simulator~\citep{goldfeder2023lightweight,goldfeder2025suite}.

We take the position that the right first entry in such a benchmark is a
\emph{structure-aware rule-based controller}: one that reads the plant topology from the
metadata and otherwise encodes physics every generated building shares. It transfers by
construction and shows which part of the score is reachable without learning.

\paragraph{Contributions.}
(i) An analysis of the SB26 reward and plant that yields the design facts above
(Section~\ref{sec:env}). (ii) SADR, a deterministic numpy controller (\PolicyMsOurs{}\,ms per
step) built from them (Section~\ref{sec:method}). (iii) A protocol of screens, loss decompositions and
pre-registered full-length confirmations on scenarios no decision touched, with Kaggle submissions used to find
which local set predicts the hidden scenarios (Section~\ref{sec:results}). (iv) An account of where the loss sits
and where it still loses (Section~\ref{sec:analysis}), every number generated from the result files.

\section{Environment and metric}
\label{sec:env}

\paragraph{Plant.}
Air handlers (AHUs) serve up to ten zones each through one VAV box per zone, mix return with
outside air (recirculation $r\le 0.7$) and can only cool: the supply is $\min(\text{mixed},
\text{SAT})$. Zone flow is $d\,F_{\max}\sqrt{P/P_{\text{design}}}$ for damper $d\in[0.1,1]$ and static
pressure $P\in[200,2000]$\,Pa; fan power is $\text{flow}\times P/\eta$, so it grows as $P^{3/2}$ once
dampers are open; the chiller's COP is a constant 3.5. Heat comes only from the VAV reheat coils on one
gas boiler loop (supply 310--350\,K, 0--20\,bar), whose effectiveness $\min(1, C_w/C_a,
(308.15-T_{\text{sa}})/(T_w - T_{\text{sa}}))$ caps the discharge at 308.15\,K; gas is billed as heat
delivered plus a tank loss.

\paragraph{Reward and metric.}
Each 5-minute step pays $0.6\,\text{(comfort shortfall)} + 0.2\,\text{(normalized energy
cost)} + 0.2\,\text{(normalized carbon)}$. The comfort shortfall is the occupancy-weighted
productivity loss (a deviation of $\Delta$\,K from the band costs $\Delta/5$ up to 5\,K and
$1+((\Delta-5)/5)^2$ beyond, normalized by total occupancy), so it is \emph{zero whenever
nobody is in the building}; occupants are present only on weekday comfort hours (arrivals
06--12, departures 13--19, none on US holidays) and only the building-wide count is
observed; the band is 294--297\,K in comfort hours and 289--298\,K otherwise. Energy is
normalized by a plant maximum that scales with $K$; in metric units one kW of electricity
costs about $2.1\times$ one kW of gas and one kW of chiller cooling about $0.6\times$. The
leaderboard metric is the mean over scenarios of
$|\text{mean per-step reward}|$, lower is better; it is dominated by the extreme climates
(the organizers' \texttt{pid\_water} loses \ExtremeRatioFairbanks$\times$ its mild-climate
mean in Fairbanks).

\paragraph{Baselines.}
The kit ships \texttt{bang\_bang}, \texttt{fixed\_setpoint}, \texttt{pid} (per-AHU PID on the
mean served-zone temperature, plant idle outside work hours) and \texttt{pid\_water}
(\texttt{pid} plus worst-zone tracking, a proportional economizer, proportional reheat and a
hot-water reset), the strongest and our main comparison; there is no learned baseline.

\section{Method: SADR}
\label{sec:method}

SADR is a single \texttt{TransferPolicy} class whose only state is per-zone integrals,
per-AHU dwell counters and the optimum-start statistics; all gains live in one parameter
dictionary (Table~\ref{tab:params}, Appendix~\ref{app:tuning}); Algorithm~\ref{alg:sadr} (Appendix~\ref{app:figs}) summarizes one step.

\paragraph{Topology from metadata.}
Index tables come from the observation field names, the positional AHU action lists, the
per-VAV action indices and the AHU-to-zone assignments; ranges and statistics from the
normalization metadata. Every action slot is written and clipped on every step; nothing depends
on the number of zones or handlers, the floor plan or the strip; non-finite inputs fall back to
schedule values; no files, no network, no randomness.

\paragraph{Modes.}
\textsc{Occupied} when the \texttt{comfort\_mode\_now} flag is set and the building is not
detected as closed; \textsc{Precondition} when unoccupied but comfort hours start within a
lead time (or \texttt{comfort\_mode\_soon}); \textsc{Setback} otherwise. ``Somebody is here''
is detected relative to the running minimum of the affinely normalized occupancy count; a
weekday on which nobody has arrived by \ParHolidayCheckH{}:00 is a holiday in the kit's
occupancy model and is treated as unoccupied until the count rises. The optimum-start lead is $\ell = \ell_0 +
\text{deficit}/(\hat{r}/s)$: the worst zone's distance from the comfort band over a
recovery-rate estimate $\hat{r}$ (K/h) re-estimated from each morning's recovery
(exponential average, weight \ParRateEma) with safety factor $s=\ParLeadSafety$, clipped to
$[\ParLeadMinH, \ParLeadMaxH]$\,h.

\paragraph{Zone loops on net heat.}
Each zone has heating and cooling targets, the sensor band shrunk by \ParMarginOcc\,K
(occupied) or \ParMarginPre\,K (preconditioning); two PI loops compute requested net heat
$q_h = k_q (T_h^\star - T)^+ + I_h$ and cooling $q_c = k_q (T - T_c^\star)^+ + I_c$ in watts
($k_q = \ParKqWPerK$\,W/K, integral gain \ParKiWPerKStep\,W/K per step); the integrals unwind
past the target, so a zone inside the band free-floats, and the larger request wins. Net
heat rather than flow because the zone-air node responds within minutes: a flow-proportional
controller with a 5-minute step oscillates, whereas a net-heat gain of the order of the air
node's conductance to the slab is stable and the integral carries the steady load.
\textsc{Setback} uses \emph{bursts} instead: a running handler pays the ventilation load of
every zone it serves, so the floor (eco band widened by \ParNightLowExt{}/\ParNightHighExt\,K)
is held by requesting \ParQSbW\,W from any zone that crosses it until it is \ParSbHystK\,K
inside, and the handler is off otherwise (holding the floor with the PI loops instead kept a
handler running for \VthreeAhuOnUnocc{} of the unoccupied steps on the coldest scenario).

\paragraph{From heat to actuators.}
A cooling zone needs capacity rate $F = q_c/(T - T_{\text{sa}})$, $T_{\text{sa}}$ being the
AHU's estimated supply temperature; a heating zone is bounded by the discharge limit, $F =
q_h/(308.15 - \delta - T)$. \emph{SAT:} a cooling handler commands the coldest value,
\ParSatCoolHi\,K: the coil load is $q_c + F\,(T_{\text{mixed}} - T)$ at a constant COP, so ventilation
term and fan power fall with flow (a Guideline-36 upward reset~\citep{ashrae2021g36,taylor2018vav} is
right in a real building, wrong here). A handler
that also heats is floored at \ParSatMixFloor\,K so its cooling zones are not overcooled into
reheat; a heating-only handler commands 305\,K so the coil idles. \emph{Static pressure} is
trim-and-respond in closed form, $P = P_{\text{design}}\max_z f_z^2$ clipped to
$[200, 2000]$\,Pa with $f_z$ each zone's requested fraction of design flow, so the most
demanding damper is fully open; dampers are $f_z/\sqrt{P/P_{\text{design}}}$. \emph{Reheat:} the discharge temperature that delivers $q_h$
at the delivered flow gives the water capacity rate $W_z = F_z\,(T_{\text{disch}} -
T_{\text{sa}})/(T_w - T_{\text{sa}})$; the loop differential pressure is reset so the largest
$W_z$ sits at \ParValveTarget{} of valve travel and each valve is $W_z/(k_v \sqrt{\Delta p})$.
An AHU runs iff a served zone requests more than \texttt{q\_on} (dwell \ParMinOnSteps{}
steps); the pump runs iff a valve is open. Recirculation is at the 0.7 cap except when a
cooling-only handler sees outside air colder than return air by more than the exhaust-fan
break-even ($0.00322$\,K/Pa$\times P$), when the outside-air fraction is solved so the mixed
air lands on the SAT target. The hot-water setpoint is 350\,K whenever heating was requested
in the last \ParHeatSeasonSteps{} steps and 310\,K otherwise: gas is billed as heat delivered,
so hotter water is strictly more useful at a tank dissipation of \TankLossExtraKw\,kW.

\paragraph{Scaling with the building (v6).}
Every VAV box in the generator has the same design flow (2\,m$^3$/s) and hot-water $k_v$, so
a 450\,m$^2$ zone has a fifth of the heating authority per square metre of a 100\,m$^2$ zone
and several times its thermal mass; v5's fixed 12\,kW cap, gains and building-wide lead were
tuned on the public zones (median \PubZoneAreaMedian\,m$^2$). v6 (i) reads each zone's floor
area from the metadata zone map (4-connected free cells in raster order reproduce the
simulator's room numbering) and scales $k_q$, $k_i$, the request cap and the setback burst with
area/\ParAreaRef\,m$^2$, clipped to $[\ParAreaScaleMin, \ParAreaScaleMax]$; (ii) bounds each
request by what the box delivers at a flow fraction $f_{\text{cap}}$ (\VSixFCapPre{} while
preconditioning or set back, \ParFCapOcc{} occupied) at the discharge limit, at 285\,K supply,
and on the water side (valve fully open at 20\,bar: asking for more raises the airflow, cools
the discharge and delivers \emph{less}); (iii) runs the optimum start per zone; (iv) raises the
setback floor of a zone whose learned rate could not recover it within the maximum lead. With
these off the code reproduces v5 bit for bit.

\paragraph{Round 4 (v7).}
\label{sec:v7}
Five switches of the v6 code (all off reproduces v6 bit for bit): (i) discharge margin $\delta$ from
\VSixDischMargin{} to \ParRFDischMargin\,K (at the cap the gas per joule of net heat is
$1 + 0.3\,(T - T_{\text{out}})/(308.15 - \delta - T)$); (ii) occupied flow above the 200\,Pa floor only in
proportion to the zone's error (from \ParRFOccEzero\,K, full \ParRFOccSpan\,K later); (iii) a night flush with
outside air at least \ParRFFlushDt\,K colder, outside heating season; (iv) in heating season the forced minimum
flow of set-back zones on a running handler is reheated to the zone temperature (1\,J of gas per joule instead
of the ventilated recovery's 1.9 later); (v) the preconditioning cap at the floor-pressure flow.

\paragraph{Tuning protocol.}
v4 came from a seeded random search (\NTuneConfigs{} configurations on six public scenarios at 1,440 steps:
defaults \TuneDefaultScreen, best \TuneBestScreen, \texttt{pid\_water} \TunePidWaterScreen; a one-at-a-time
ablation shows that doubling the integral gain to \ParKiWPerKStep\,W/K per step is the whole gain,
\AblRelDefAddKi), v5 from \NCandV{} physics-motivated candidates run at full length on the public set
(Appendix~\ref{app:tuning}). Rounds 3--6 screened single changes on public and held-out scenarios (2,880
steps; round 6 at full length), combined the helpful ones and shipped only under a rule fixed beforehand (not
worse on the public set, better on held-out scenarios at full length); rounds 4--6 wrote the protocol down
before running anything (\texttt{PREREG.md}) and confirm on scenarios no decision touched: round 3 on
\NHeldAll{} scenarios of \NHeldAllBuild{} default-range buildings (\NHeldUnscr{} never screened), round 4 on
\RFFreshNBuild{} freshly generated buildings (three candidates added after a first look at the screen), round 5
on the shipped six under the one station the kit ships but no public scenario uses (Minneapolis), round 6 on the
shipped six under the harshest strips available (Section~\ref{sec:results}).

\section{Results}
\label{sec:results}

\paragraph{Setup.}
All runs use the kit's environment loader and \texttt{SimulationTracker}, 8,640 steps per scenario,
and the kit's \texttt{score\_leaderboard.py} formula with paired bootstrap CIs over scenarios. Rounds 1--4
ran on an Apple M5 Pro and a Windows laptop, round 5 on the laptop, round 6 on private Kaggle CPU kernels and the
M5 Pro; the seeds are fixed and runs reproduce across machines (SADR bit for bit between laptop and Kaggle,
\RSixReproIdentical{} of \RSixReproN{} checks; within \RSixMacReproMaxRel{} on the M5 Pro in round 6; \texttt{pid\_water}
within \MachineGapPidWater). Local \texttt{pid\_water} returns differ
from the shipped \texttt{baseline\_returns.json} by \GapMeanPidWater{} on average, so all comparisons use
local runs.

\begin{table}[t]
\caption{Leaderboard metric (mean over the \NScen{} public scenarios of $|$mean per-step
reward$|$; lower is better) and per-station means. Paired bootstrap 95\% CIs (\NBoot{}
resamples) are for SADR v6 minus the row (v7 against v6: Table~\ref{tab:r4}; v8 against v7 and v9 against v8:
Section~\ref{sec:results}). v9 ran on the M5 Pro, v8 on the laptop, the other rows on one machine (reproducible across
machines); the last row is the organizers' shipped returns. Per-scenario numbers: Appendix~\ref{app:per_scenario}.}
\label{tab:main}
\centering\scriptsize
\setlength{\tabcolsep}{3pt}
\begin{tabular}{lcccccc c}
\toprule
policy & all & Fairbanks & Denver & Atlanta & Phoenix & Singapore & $\Delta$ vs v6 [95\% CI] \\
\midrule
SADR v9 (final, round 6) & \textbf{\RSixConfPubCand} & \RSixConfPubCandFairbanks & \RSixConfPubCandDenver & \RSixConfPubCandAtlanta & \RSixConfPubCandPhoenix & \RSixConfPubCandSingapore & Sec.~\ref{sec:results} \\
SADR v8 (round 5) & \RFiveConfPubCand & \RFiveConfPubCandFairbanks & \RFiveConfPubCandDenver & \RFiveConfPubCandAtlanta & \RFiveConfPubCandPhoenix & \RFiveConfPubCandSingapore & Sec.~\ref{sec:results} \\
SADR v7 (round 4) & \RFConfPubCand &\RFConfPubCandFairbanks & \RFConfPubCandDenver & \RFConfPubCandAtlanta & \RFConfPubCandPhoenix & \RFConfPubCandSingapore & Table~\ref{tab:r4} \\
SADR v6 (round 3) & \MetricOurs & \MetricOursFairbanks & \MetricOursDenver & \MetricOursAtlanta & \MetricOursPhoenix & \MetricOursSingapore & -- \\
SADR v5 & \MetricOursVFive & \MetricOursVFiveFairbanks & \MetricOursVFiveDenver & \MetricOursVFiveAtlanta & \MetricOursVFivePhoenix & \MetricOursVFiveSingapore & \DiffOursVFive{} [\DiffLoOursVFive, \DiffHiOursVFive] \\
pid\_water & \MetricPidWater & \MetricPidWaterFairbanks & \MetricPidWaterDenver & \MetricPidWaterAtlanta & \MetricPidWaterPhoenix & \MetricPidWaterSingapore & \DiffPidWater{} [\DiffLoPidWater, \DiffHiPidWater] \\
fixed\_setpoint & \MetricFixedSp & \MetricFixedSpFairbanks & \MetricFixedSpDenver & \MetricFixedSpAtlanta & \MetricFixedSpPhoenix & \MetricFixedSpSingapore & \DiffFixedSp{} [\DiffLoFixedSp, \DiffHiFixedSp] \\
bang\_bang & \MetricBangBang & \MetricBangBangFairbanks & \MetricBangBangDenver & \MetricBangBangAtlanta & \MetricBangBangPhoenix & \MetricBangBangSingapore & \DiffBangBang{} [\DiffLoBangBang, \DiffHiBangBang] \\
pid\_water (organizer values) & \MetricPidWaterOff & \MetricPidWaterOffFairbanks & \MetricPidWaterOffDenver & \MetricPidWaterOffAtlanta & \MetricPidWaterOffPhoenix & \MetricPidWaterOffSingapore & \DiffPidWaterOff{} [\DiffLoPidWaterOff, \DiffHiPidWaterOff] \\
\bottomrule
\end{tabular}
\end{table}

\paragraph{Main result.}
SADR v9 reaches \RSixConfPubCand{} on the public set, v8 \RFiveConfPubCand{} and v7 \RFConfPubCand, against \MetricPidWater{} for
\texttt{pid\_water} (Table~\ref{tab:main}; v7's ratio \RFPwPubRatio, 95\% CI \RFPwPubRatioLo{} to \RFPwPubRatioHi, wins
\RFPwPubWins). On the public set v4 to v5 cut the metric by \CandWinnerRelRed{} and v5 to v6 by \RelRedOursVFive{} (CI
\RelRedLoOursVFive{} to \RelRedHiOursVFive), small because the public zones are small.

\paragraph{Where the loss goes.}
Per-step plant energy logs split the loss (Figure~\ref{fig:decomp}): on the public set v5 pays
\DecPubVFiveGasShare{} for gas, \DecPubVFiveCoolingShare{} for the chiller, \DecPubVFiveFanShare{} for fans and
\DecPubVFiveComfortShare{} for comfort (Fairbanks: gas \DecPubVFiveGasShareFairbanks); \texttt{pid\_water} pays
\RegretSharePidWater{} for comfort.

\paragraph{Held-out buildings.}
On \NHeldAll{} scenarios of \NHeldAllBuild{} default-range generated buildings (Figure~\ref{fig:heldout}) v6 scores
\HeldAllMetricOurs{} against \HeldAllMetricOursVFive{} for v5 and \HeldAllMetricPidWater{} for \texttt{pid\_water}
(\HeldUnscrVFiveRelRed{} below v5 on the \NHeldUnscr{} never screened; Table~\ref{tab:candidates_r3}).

\paragraph{Round 4.}
Of \RFNSingles{} single changes (Table~\ref{tab:candidates_r4}) \RFNHelps{} helped and \RFNWorse{} were worse (caps
below the floor-pressure flow; the night flush in heating season). The composite v7 beats v6 at full length on every
set (Table~\ref{tab:r4}): \RFConfPubRed{} on the public set (CI \RFConfPubRedLo{} to \RFConfPubRedHi),
\RFConfOldCleanRed{} and \RFConfFreshRed{} on the clean old and fresh held-out sets, mostly gas.

\input{round6_kaggle_table}

\paragraph{Kaggle.}
On the Kaggle public leaderboard (about ten of 30 hidden scenarios: six buildings from the same generator under
five stations absent from the public set) v7 scores \RFiveKagVSeven{} (rank \RFiveKagRankVSeven/\RFiveKagTeamsVSeven),
\RFiveKagGapPub{} times its public-set number, although Kaggle's run of a public scenario matches the local total to
0.02\%. Round 5 used Kaggle as an oracle: \RFiveNProbes{} submissions of fixed files (v6, v5, \texttt{pid\_water})
and a pre-registered comparison of each local set's level and shape $E$, how well it reproduces each policy's ratio
to v7 (Table~\ref{tab:kaggle_r5}). v5's ratio to v7 on Kaggle (\RFiveKagRatioVFive) is the public set's
(\RFivePubRatioVFive), not that of the larger default-range buildings (\RFiveOldRatioVFive): the hidden buildings
behave like the shipped six. \texttt{pid\_water}'s ratio (\RFiveKagRatioPid) falls in cold weather, where both
controllers pay the same ventilation physics (public \RFivePubRatioPid, Minneapolis \RFiveMinnRatioPid), and the
public set plus Minneapolis reproduces it (\RFivePubMinnRatioPid; $E=\RFivePubMinnShapeErr$ against
\RFivePubShapeErr): the gap is weather, and as v7 reaches only \RFivePubMinnLevel{} of its Kaggle score there, the
hidden stations are harsher still. A public notebook's \texttt{pid\_water} scores \RFiveThirdAKag{} on Kaggle (ours
\RFiveKagPid). Round 5's rule required its candidate to win on P+M and on the fresh default-range set FD.

\paragraph{Round 5.}
Of \RFiveNSingles{} single changes to v7 screened at \RFiveScreenSteps{} steps (Table~\ref{tab:candidates_r5}),
\RFiveNHelps{} helped, \RFiveNNs{} were not significant and \RFiveNWorse{} were worse (a longer lead before Mondays,
\RFiveMonHeldRelAbs; recovering at the flow cap from the start, \RFivePreburstHeldRelAbs). The composite v8 of a deeper
setback far from the next workday (floor 6\,K lower, ceiling 4\,K higher, ramped out 10--22\,h before the start) and a
request cap of 20\,kW per 100\,m$^2$ (it bound at half of what box and coil deliver in deep cold) lowers the loss at
full length by \RFiveConfPubRelAbs{} on the public set (CI \RFiveConfPubRelHiAbs{} to \RFiveConfPubRelLoAbs; wins
\RFiveConfPubWins), \RFiveConfMinnRelAbs{} on M and \RFiveConfPubMinnRelAbs{} on P+M (CI \RFiveConfPubMinnRelHiAbs{} to
\RFiveConfPubMinnRelLoAbs) and ties on FD (\RFiveConfFreshDefRelAbs; the CI's upper end is below zero for
\RFiveSeedsFreshDefBelow{} of 10 bootstrap seeds). It met the rule as written; Kaggle: \RFiveKagVEight{} against v7's
\RFiveKagVSeven, within the noise of the public split.

\paragraph{Round 6.}
Round 6 followed the probes to harsh weather with two new sets of the shipped six: a confirmation set HC
(\RSixHarshCN{} scenarios: the Minneapolis strips not used before, the coldest Fairbanks and Denver strips and new
30-day strips cut from the coldest months of the shipped weather files) and a screening set HS (\RSixHarshSN). v8
scores \RSixHarshCVEight{} on HC, \RSixHarshCLevel{} times its Kaggle score (public set \RSixPubLevel;
Table~\ref{tab:kaggle}); a mixture of \RSixMixWPct{} HC and the rest public, fitted to v8's Kaggle score, predicts
\texttt{pid\_water}'s Kaggle ratio (\RSixMixPid{} against \RSixKagRatioPid). The v8 diagnosis
(Section~\ref{sec:analysis}) dropped six planned changes before any ran: an outdoor-temperature term in the optimum
start (two settings) and four whose whole target is under 1\% of the HS loss or out of the controller's reach
(supply-air reset under heating load, hot-water reset, no simultaneous reheat and cooling, a lower occupied heating
margin). Of \RSixNSingles{} single changes screened at full length on HS and twelve public scenarios
(Table~\ref{tab:candidates_r6}; one added after a first look), \RSixNHelps{} helped and one did nothing. The largest gains came from
reheating the forced minimum air of set-back zones on running handlers to the discharge limit instead of to the zone
temperature, which buys heat at 1\,J of gas per joule against about 1.9 in the zone's own ventilated recovery (before a
workday \RSixScrCElevenHsRel{} on HS, CI \RSixScrCElevenHsLo{} to \RSixScrCElevenHsHi; near the setback floor
\RSixScrCNineHsRel), and from a deeper weeknight setback (\RSixScrCEightHsRel); a nearer preconditioning target
(\RSixScrCSixHsRel) and shorter leads (\RSixScrCSevenHsRel) helped less; joining another zone's burst did nothing. A
5$^\circ$C frost floor, run only to price it, raises the HS loss by \RSixScrCTenHsRelAbs. The composite of the best change per
family (\RSixCompMembers; v9) lowers HS by \RSixScrCompHsRelAbs{} (CI \RSixScrCompHsHiAbs{} to \RSixScrCompHsLoAbs) and the
public twelve by \RSixScrCompPubRelAbs{} and is chosen. At full length v9 lowers the loss on HC by \RSixConfHarshCRelAbs{} (CI
\RSixConfHarshCRelHiAbs{} to \RSixConfHarshCRelLoAbs; wins \RSixConfHarshCWins; upper end below zero for
\RSixSeedsHarshCBelow{} of 10 bootstrap seeds) and on M by \RSixConfMinnRelAbs, and changes the public set by
\RSixConfPubRel{} (CI \RSixConfPubRelLo{} to \RSixConfPubRelHi): the rule (public not worse, HC better) is met and v9 is
the submission. It does not help on FD (\RSixConfFreshDefRel, CI \RSixConfFreshDefRelLo{} to \RSixConfFreshDefRelHi; wins
\RSixConfFreshDefWins; comfort rises by more than the total) and is up to \RSixConfPubWorstRel{} on hot public strips,
which the two hot screening scenarios under-sampled; the ledger points to the nearer preconditioning target (c06): with the shorter lead in its
place, a post-hoc composite (ledger only) changes the public set by \RSixCompBPubRel{} (CI \RSixCompBPubRelLo{} to
\RSixCompBPubRelHi; Table~\ref{tab:candidates_r6}). On Kaggle c01 scores \RSixProbeAScore{} (v8 \RSixKagVEight, within
noise); the v9 probe was pending.

\section{Analysis}
\label{sec:analysis}

\paragraph{Why it wins, and where the loss sits.}
\texttt{pid\_water} keeps \OccOutPidWater{} of occupied zone-steps outside the band (one loop per handler,
a step start at 06:00; Figure~\ref{fig:hours}); SADR v6 keeps \OccOutOurs{} and uses far less energy
(floor pressure with wide dampers, least reheat flow at the discharge cap, no conditioning while nobody
is charged for comfort). For v6, \UnoccShareOurs{} of the loss is incurred on unoccupied steps, most at
hour \PeakLossHourOurs{}, the warm-up before 06:00. Fairbanks is gas: with a mandatory 30\% outside-air
fraction at $-20^\circ$C heat costs about $1.9$\,J of gas per joule at the discharge cap; Phoenix and
Singapore are chiller electricity, where the coldest supply at the lowest flow is cheapest.

\paragraph{What rounds 3 and 4 found.}
v6 cut the largest-quartile recovery of v5's 350--430\,m$^2$ zones from \DecHeldRVFiveRecoveryBig\,h to
\DecHeldROursRecoveryBig\,h. The generated buildings' larger chiller term is a normalization effect: they cool
\RFCoolWmRThree{} against \RFCoolWmPub\,W/m$^2$ but have \RFAreaPerAhuRThree{} instead of \RFAreaPerAhuPub\,m$^2$ of zone per handler.

\paragraph{What round 5 found.}
Under the Minneapolis strips (set M) v7 pays \RFiveDecMinnGasShare{} of its loss for gas and
\RFiveDecMinnComfortShare{} for comfort, \RFiveCapComfortShareOfMinn{} of it in one January scenario whose
\RFiveCapNZones{} worst zones sat \RFiveCapDefMin--\RFiveCapDefMax\,K below the band because the request cap
(\RFiveCapReqKwMin--\RFiveCapReqKwMax\,kW) bound at about half of what box and coil deliver (v8 lifts it). Handlers
start \RFiveAhuStartsMin--\RFiveAhuStartsMax{} times a day, so cycling is not a lever. Starter-kit pull request~\#2
(unmerged) bills boiler gas from a tank energy balance and closes a gap that rewards saw-toothing the hot-water
setpoint (heat billed at 0.38 of its value); in a patched kit v7's loss falls by \RFivePrTwoVSevenRelAbs{} on the
\RFivePrTwoVSevenN{} public Fairbanks scenarios while \texttt{pid\_water}'s rises by \RFivePrTwoPidRelAbs: SADR's score
does not rest on that gap.

\paragraph{What round 6 found.}
On the harsh screening set v8 pays \RSixDecGas{} of its loss for gas: \RSixDecOccWd{} of the gas falls in
occupied weekday hours, \RSixDecPreWd{} in weekday preconditioning (05:00--06:00 alone \RSixDecHourFive) and
\RSixDecSbWe{} in weekend setback. The outside-air load of running handlers equals \RSixDecVent{} of the gas,
\RSixDecVentHeated{} of it through zones that are being heated: the handler cannot heat, the mixed air is at
least 30\% outside air, and at the discharge cap that overhead is fixed per joule, so in occupied hours it is
physics, not control. The morning recovery rate does not fall in colder weather (within a zone it rises with the
lift, \RSixRecSlope{}\,K/h per K, CI \RSixRecSlopeLo{} to \RSixRecSlopeHi, \RSixRecN{} zone-mornings), so an
outdoor-temperature term in the optimum start had nothing to correct. v8's weekend floor is about 2$^\circ$C, and
burst cycling at $-38^\circ$C lets the coldest zone reach \RSixDecColdestWe$^\circ$C (\RSixDecBelowFive{} of
zone-steps below 5$^\circ$C): the simulator does not charge for frost; a real building would.

\paragraph{Where it loses.}
SADR v7 is worse than \texttt{pid\_water} on \RFNScenWorseCand{} of \NScen{} public scenarios;
its largest residual loss is on \RFMaxCandScen{} (\RFMaxCandScore), almost all gas. What remains is the
morning recovery in the coldest strips, kept cheap rather than perfect (Figure~\ref{fig:hours}), and box
capacity on the largest zones; a learned policy's advantage should be looked for in the recovery.

\section{Limitations}
SADR encodes simulator constants (discharge limit, valve $k_v$, constant COP) that a real building would
commission, and its weekend floor near freezing ignores frost protection. Selection used the public set and
screens, so the public number is optimistic. MPC or a learned residual on top of SADR is the next step.

\section{Reproducibility and disclosure}
Code: \url{https://github.com/neelsirivara/neurips2026-smart-buildings}; \texttt{submission/transfer\_policy.py} is
the policy, \texttt{experiments/src/} generates every number (\texttt{finalize.sh}); seeds fixed; CPU only. An LLM
coding assistant (Claude) wrote code and drafted this text; every number is checked in \texttt{paper/AUDIT.md}.

\bibliographystyle{plainnat}
\bibliography{refs}

\clearpage
\appendix
\section{Per-scenario results}
\label{app:per_scenario}
Tables~\ref{tab:per_scenario}, \ref{tab:per_scenario_held} and~\ref{tab:per_scenario_held_r3}
list every scenario's total return and leaderboard score (SADR v6), and
Table~\ref{tab:per_scenario_fresh} the round-4 confirmation set, Table~\ref{tab:r4} round 4 at full length and
Table~\ref{tab:per_scenario_r5m} the round-5 Minneapolis set; all values are generated from the
result files by \texttt{make\_macros.py}.
\begin{table}[h]
\caption{Round 4 at full length (8,640 steps): the composite v7 (\RFChosen) against the shipped
v6 on the same scenarios; $\Delta$ = paired mean difference in $10^{-5}$ per step with its
bootstrap 95\% CI (\NBoot{} resamples; negative favours v7). Old clean = the round-2/3 held-out
scenarios no screen touched; fresh = \RFFreshNBuild{} buildings generated for round 4 and run only
for this table, half in the shipped six's parameter ranges, half in the kit's default ranges.}
\label{tab:r4}
\centering\scriptsize
\setlength{\tabcolsep}{4pt}
\begin{tabular}{lrrrrrr}
\toprule
set & $n$ & v6 & v7 & $\Delta$ [95\% CI] & relative [95\% CI] & wins \\
\midrule
public & \RFConfPubN & \RFConfPubVSix & \RFConfPubCand & \RFConfPubDiff{} [\RFConfPubLo, \RFConfPubHi] & \RFConfPubRel{} [\RFConfPubRelLo, \RFConfPubRelHi] & \RFConfPubWins \\
old held-out, clean & \RFConfOldCleanN & \RFConfOldCleanVSix & \RFConfOldCleanCand & \RFConfOldCleanDiff{} [\RFConfOldCleanLo, \RFConfOldCleanHi] & \RFConfOldCleanRel{} [\RFConfOldCleanRelLo, \RFConfOldCleanRelHi] & \RFConfOldCleanWins \\
fresh, shipped ranges & \RFConfFreshShipN & \RFConfFreshShipVSix & \RFConfFreshShipCand & \RFConfFreshShipDiff{} [\RFConfFreshShipLo, \RFConfFreshShipHi] & \RFConfFreshShipRel{} [\RFConfFreshShipRelLo, \RFConfFreshShipRelHi] & \RFConfFreshShipWins \\
fresh, default ranges & \RFConfFreshDefN & \RFConfFreshDefVSix & \RFConfFreshDefCand & \RFConfFreshDefDiff{} [\RFConfFreshDefLo, \RFConfFreshDefHi] & \RFConfFreshDefRel{} [\RFConfFreshDefRelLo, \RFConfFreshDefRelHi] & \RFConfFreshDefWins \\
fresh, all & \RFConfFreshN & \RFConfFreshVSix & \RFConfFreshCand & \RFConfFreshDiff{} [\RFConfFreshLo, \RFConfFreshHi] & \RFConfFreshRel{} [\RFConfFreshRelLo, \RFConfFreshRelHi] & \RFConfFreshWins \\
clean held-out, pooled & \RFConfPoolN & \RFConfPoolVSix & \RFConfPoolCand & \RFConfPoolDiff{} [\RFConfPoolLo, \RFConfPoolHi] & \RFConfPoolRel{} [\RFConfPoolRelLo, \RFConfPoolRelHi] & \RFConfPoolWins \\
\bottomrule
\end{tabular}
\end{table}
\input{per_scenario_table}
\input{per_scenario_table_r4}
\input{per_scenario_table_r5}
\begin{figure}[h]
\centering
\figfile{0.92\linewidth}{figures/fig_round4_buildings.pdf}
\caption{Round 4: change of the metric from v6 to v7 on every generated held-out building
(diamonds: the building's five strips pooled; dots: single scenarios; negative is better), each
run once at full length; buildings grouped by set and ordered by zone count.}
\label{fig:r4buildings}
\end{figure}
\begin{figure}[h]
\centering
\figfile{0.86\linewidth}{figures/fig_scores_by_scenario.pdf}
\caption{Per-scenario leaderboard score (log scale, lower is better), SADR v6 vs
\texttt{pid\_water} on the same machine, grouped by station and labelled by building (b =
public, g = generated held-out). Left: the 30 public scenarios; right: the \NHeld{}
held-out scenarios of buildings 0501 and 0502.}
\label{fig:scores}
\end{figure}

\section{Parameters, ablation and candidates}
\label{app:tuning}
\begin{table}[h]
\caption{Final SADR parameters (v9); values that changed from v4 to v5 are marked; the fifth row
is new in v6 (round 3), the discharge margin (v6: \VSixDischMargin\,K), the preconditioning cap
(v6: \VSixFCapPre) and the seventh row are round 4, the eighth row round 5, the last row and the precondition margin
and setback widening (v8: 0.8\,K and 8\,K) round 6. The full dictionary is in the code.}
\label{tab:params}
\centering\footnotesize
\begin{tabular}{llll}
\toprule
occupied margin & \ParMarginOcc\,K & precondition margin & \ParMarginPre\,K \\
setback widening (low/high) & \ParNightLowExt{}/\ParNightHighExt\,K (v4: \VFourNightLowExt/\VFourNightHighExt) & $k_q$, $k_i$ & \ParKqWPerK, \ParKiWPerKStep\,W/K \\
cooling SAT / mixed floor & \ParSatCoolHi{}/\ParSatMixFloor\,K (v4 ceiling: \VFourSatCoolHi) & discharge margin & \ParDischMarginK\,K \\
lead safety / rate$_0$ & \ParLeadSafety{} / \ParRateZeroKph\,K/h & lead range & \ParLeadMinH--\ParLeadMaxH\,h \\
heating-season memory & \ParHeatSeasonSteps{} steps & AHU dwell & \ParMinOnSteps{} steps \\
area scaling (v6) & per \ParAreaRef\,m$^2$, clip $[\ParAreaScaleMin, \ParAreaScaleMax]$ & flow cap $f_{\text{cap}}$ (v6) & \ParFCapPre{} pre/setback, \ParFCapOcc{} occ. \\
occupied pressure cap (v7) & from \ParRFOccEzero\,K, span \ParRFOccSpan\,K & night flush (v7) & $\ge$\ParRFFlushDt\,K, $\le$\ParRFFlushWindow\,h before work \\
weekend setback (v8) & $-$\RFiveParWkLo/$+$\RFiveParWkHi\,K, out \RFiveParWkRampEnd--\RFiveParWkRampStart\,h before work & request cap (v8) & \RFiveParQcapKw\,kW/100\,m$^2$ (v7: \RFiveParQcapKwVSeven) \\
set-back air boost (v9) & $\le$8\,h before work, outside $\le$ floor & near-floor boost (v9) & within 2\,K of the floor \\
\bottomrule
\end{tabular}
\end{table}

Table~\ref{tab:ablation} reverts each tuned v4 value on its own (and adds each on its own to
the defaults) on the screening set; Table~\ref{tab:candidates} lists every v5 candidate at
full length on the 30 public scenarios against v4 on the same machine;
Table~\ref{tab:candidates_r3} lists every round-3 (v6) variant on both screening sets against
v5, Table~\ref{tab:candidates_r4} every round-4 candidate against v6 and Table~\ref{tab:candidates_r5}
every round-5 candidate against v7, Table~\ref{tab:kaggle_r5} the round-5 comparison of local sets with Kaggle
(v7 reference) and Table~\ref{tab:candidates_r6} every round-6 candidate against v8. The ideas tried and rejected,
with their numbers, are in \texttt{experiments/results/round3/candidates.json}, \texttt{round4/candidates.json},
\texttt{round5/screen\_summary.json} and \texttt{round6/screen\_summary.json}.
\input{tuning_tables}
\input{tuning_tables_r5}
\input{round5_kaggle_table}
\input{tuning_tables_r6}

\section{Algorithm and additional figures}
\label{app:figs}
\begin{figure}[h]
\centering
\begin{minipage}[t]{0.5\linewidth}\centering
\figfile{\linewidth}{figures/fig_loss_decomp.pdf}
\caption{Per-station decomposition of the per-step loss (metric units) into comfort, fan,
chiller, pump and gas on the 30 public scenarios, SADR v5 and v6.}
\label{fig:decomp}
\end{minipage}\hfill
\begin{minipage}[t]{0.44\linewidth}\centering
\figfile{\linewidth}{figures/fig_heldout_buildings.pdf}
\caption{Metric per generated held-out building (log scale, lower is better; ordered by zone
count): \texttt{pid\_water}, SADR v5 and v6, each run once at full length on the building's
five strips.}
\label{fig:heldout}
\end{minipage}
\end{figure}
\begin{algorithm}[h]
\caption{SADR, one step (vectorized over zones $z$ and handlers $k$)}
\label{alg:sadr}
\begin{algorithmic}[1]
\State mode $\gets$ \textsc{Occupied} / \textsc{Precondition} / \textsc{Setback} from flags, time, occupancy count, lead $\ell$
\State $(T_h^\star, T_c^\star)_z \gets$ band for the mode; update $I_h, I_c$; $q_h, q_c \gets$ PI requests (W), or setback bursts
\State run$_k \gets$ hysteresis on $\max_{z\in k}(q_h, q_c)$
\State SAT$_k \gets$ \ParSatCoolHi\,K if cooling (\ParSatMixFloor\,K if also heating), 305\,K otherwise; recirc$_k \gets$ economizer
\State $F_z \gets q_c/(T_z - T_{\text{sa},k})$ or $q_h/(308.15-\delta-T_z)$; $P_k \gets P_{\text{design}}\max_z f_z^2$; damper$_z \gets f_z/\sqrt{P_k/P_{\text{design}}}$
\State $W_z \gets F_z (T_{\text{disch},z} - T_{\text{sa}})/(T_w - T_{\text{sa}})$; $\Delta p \gets (\max_z W_z / (0.9 k_v))^2$; valve$_z \gets W_z/(k_v\sqrt{\Delta p})$
\State $T_w^{\text{sp}} \gets 350$\,K if heating season else 310\,K; pump $\gets$ any valve open
\State scatter all setpoints into the normalized action vectors; clip
\end{algorithmic}
\end{algorithm}
\begin{figure}[h]
\centering
\figfile{\linewidth}{figures/fig_loss_by_hour.pdf}
\caption{Mean per-step loss by hour of day over the \NAnalysis{} public scenarios, split
into comfort regret and energy\,+\,carbon; comfort hours (06--19) shaded. Left: SADR v6;
right: \texttt{pid\_water}.}
\label{fig:hours}
\end{figure}
\begin{figure}[h]
\centering
\figfile{\linewidth}{figures/fig_trace_building_train_8_fairbanks_8_day8.pdf}
\caption{One weekday on \texttt{building\_train\_8/fairbanks\_8}: zone-temperature
envelope (min--max shaded, median line) against the comfort band (green wash), with the
building occupancy count shaded in grey. Left: SADR v6; right: \texttt{pid\_water}.}
\label{fig:trace}
\end{figure}
\begin{figure}[h]
\centering
\figfile{0.5\linewidth}{figures/fig_tuning.pdf}
\caption{Screening score (six scenarios, five days) of every random-search configuration,
sorted, with the design defaults highlighted and \texttt{pid\_water} on the same runs.}
\label{fig:tuning}
\end{figure}

\end{document}
